\documentclass{article}
\usepackage[utf8]{inputenc}

\title{A Very Cool and Catchy Title}
\author{Zhenkai Gao, Alex Kautz, Pascale Leone
}
\date{\vspace{-1em}}

\usepackage{natbib}
\usepackage{graphicx}
\usepackage{parskip} % Paragraphs use vertical spacing instead of indentation.

\begin{document}

\maketitle

% for final report only
%\begin{abstract}
%A one paragraph high level description of the work. The abstract is typically the first thing someone would read from the paper before deciding whether to continue reading and hence, serves as an advertisement to the reader to read the whole paper. 
%\end{abstract}

\section{Introduction}

\textit{The introduction should describe the problem (in a non-technical manner, i.e., without math, equations, etc.), as well as motivate the problem, i.e., why is it important?}

Imagine a future where you receive a box of groceries from a delivery service (like Amazon), take the box to your kitchen table, and open it. Then, a helpful robotic arm takes items out of the box and puts them away. The arm knows the best place for each item, be it the cupboard, fridge, or freezer.

This kind of assistance can be a powerful utility for elderly or disabled individuals, as well as for professional kitchens where inventory management is a major time sink.

For our project, we simulated a version of this arm using Gazebo in order to explore the problem space. We designed an environment analogous to a real kitchen, and laid out a design for the full system.

We then focused on just the sub-problem of picking up objects, integrating computer vision techniques like YOLO, and trajectory planning tools like MoveIt.

We analyzed the effect inverted images and shadows may have on the YOLO algorithm. We also measured the effect that lowering the position and orientation tolerance in trajectory planning has on the reliability of the picking up task.

We hope this work can lead to real-world implementations of the picking up system.


\section{Background Related Work}

\textit{Following, you should provide the necessary background and discuss related work in the RL literature. This section should also be about a page. Citations should be in BibTeX format \citep{thrun2005probabilistic}.}

\textbf{We should copy the background and related work from our proposal and then expand upon it.}


\textbf{We can expand upon it by rewriting our notes from each article we read.}

Much research has been done in the field of home service robotics. For example, assisting the elderly with household tasks is a common application for robotics. In 1998, a team based out of Italy developed MOVAID. The MOVAID system consisted of a robotic arm attached to a driving base. MOVAID's goal was to assist in 3 main tasks: heating up food, cleaning the kitchen, and changing bed sheets. As part of that first task, the robot had to take food out of the refrigerator \cite{doi:10.3233/TAD-1999-10202}. 

Researchers at the School of Artificial Intelligence, Jingchu University of Technology present a novel framework for an in-home service robot \cite{11179321}. They utilize ROS to develop high-precision robotic manipulators
for enhanced grasping capabilities. They evaluate the success of their robot in a simulated home environment. In addition, the authors employ vision model YOLOV5 for object detection and recognition.

Similar to these works, we want to use a robotic arm to unpack and place groceries. We believe this task is important because it can help the elderly or disabled in the kitchen. We use a static robot arm rather than a moving base and focus on the task of putting away groceries. This task involves three key areas: object detection, trajectory planning, and stable placement.

Object detection in a cluttered environment is an important problem. Researchers at Jaume I University in Castellon de la Plana, Spain studied object detection in a cluttered scene in order to help the BAXTER robot pick the correct object in a cluttered environment \cite{10.1145/3164541.3164628}. They employed deep learning methods to improve performance, and compare using RESNET, a semantic segmentation based method, to a localization method with YOLO for object detection. 

Trajectory planning is also relevant to our problem. Our robotic arm must navigate from the box of groceries to the appropriate storage container, without hitting any obstacles. Researchers in the Department of Mechatronics Engineering at the Isparta University of Applied Sciences in Isparta Turkey studied trajectory planning for a robotic arm with 6 degrees of freedom \cite{EKREM2023106099}. Their experiments focused on the Particle Swarm Optimization (PSO) algorithm, which is a population-based stochastic optimization method. This algorithm starts with a population of random solutions and searches for the optimal solution by updating generations. They proposed that the PSO algorithm optimizes trajectory planning of the 6-degree-of-freedom. 

It is also important that our robot is able to facilitate stable placement of each grocery item. As the storage containers accumulate objects, the robot must determine an appropriate location for the new grocery item. Haustein, Hang, Stork, and Kragic address the problem of placement planning by creating an algorithm that computes placements for objects based on three main principles: (1) suitable locations must be identified, (2) placements must be reachable, and (3) not all placements are equally desirable \cite{8967732}.

Recent work in robotic manipulation has shifted from traditional model- and planning-based approaches toward data-driven and learning-based methods. Kleeberger et al. survey this transition, emphasizing the rise of learning-based grasping and its challenges in generalization \cite{kleeberger2020survey}. While our project focuses on a single robotic arm, real-world home service robots require whole-body coordination, including mobile bases and multiple manipulators. Jiang et al. propose a benchmark and system framework for real-world household tasks \cite{jiang2025behavior}.

In addition, recent approaches explore unified learning frameworks that jointly model perception, navigation, and interaction, enabling robots to actively improve performance through interaction with the environment \cite{uppal2024spin}. Finally, vision-language models allow robots to perform more semantic, task-aware manipulation, such as grasping specific object parts and generalizing to unseen objects \cite{oort2025open}.


\section{Technical Approach / Methodology / Theoretical Framework}

% for proposal
%Describe how you will approach the problem and its technical formulation. Feel free to re-state the basic RL formulas (e.g., if using Q-learning, state the update rule or the formula for what the Q function approximates). 


% for final report
\textit{A detailed description of your problem (with math, notation, algorithms, figures, etc.). Use footnotes to cite links to your code or videos. Add a footnote like "All developed source code for this project is available at ..."}

\textit{Notation and Problem Formulation
How did you formalize the problem using mathematical notation?
Tasks
Subsections are useful for breaking down the problem into sub-parts. For example, you could break down the tasks for your project and list them one by one. 
}

\subsection{Theoretical Framework}

\subsubsection{Computer Vision}

To detect, recognize, and locate our object we use a YOLO (You Only Look Once) model. YOLO is an object detection and segmentation model initially developed by Joseph Redmon, Santosh Divvala, Ross Girshick, and Ali Farhadi at the University of Washington in 2015 \cite{DBLP:journals/corr/RedmonDGF15}. Their model differs from prior object detection techniques because they frame  object detection as a regression problem. A single convolutional neural network simultaneously predicts bounding boxes and class probabilities. This allows for incredibly fast and generalizable object detection.


In our project, we use the COCO-pretrained YOLOv8 model, released in 2023 by Ultralytics \cite{yolov8_ultralytics}. COCO is a large-scale annotated object detection and segmentation dataset. The classes in this dataset are diverse, containing categories such as cars, animals, or handbags. This dataset is appropriate for the scale of our problem, as it includes several grocery items as classes, specifically apple, orange, banana, carrot, and broccoli. Since we are using fixed items *blank* this is adequate.

Our overhead camera uses YOLOv8 to detect the apples bounding box. From there it computes and stores the centroid of the apple. We then compute the depth level based on the centroid coordinates. With these x,y,z values we transform the pixel location to the location in the simulation world. This information is published for use in the trajectory calculation.
\subsubsection{Trajectory Calculation}

How we used Move-It! to calculate the trajectory while avoiding objects

\subsubsection{Destination Determination}

How we used a lookup table after identifying the object from computer vision

\subsection{Technical Approach}

\subsubsection{Simulation Environment}

For our simulation, we employed Gazebo~\cite{koenig2004gazebo}, an open-source robotics simulator with deep integration into ROS~\cite{quigley2009ros}. Gazebo enables the construction of digital replicas of physical robots within customizable 3D environments.

We designed an environment that was as simple as possible while still representing our project. We used a two-tier shelf to represent the cabinet, refrigerator, and freezer combination. This retained the concept of having different locations to place objects. We used a simple surface opposite the robot in place of the table and box that contain the groceries. Finally, for the groceries themselves, we used basic shapes, such as a sphere for the apple. You can see the environment running inside Gazebo in Figure \ref{fig:environment}.


\begin{figure}[h]
    \centering
    \includegraphics[width=0.7\textwidth]{documents/latex/final_report/Env.png}
    \caption{Pictured in the environment: (1) two tier shelf, left; (2) robotic arm, center; (3) overhead camera, center; (4) table, bottom right; (5) apple, bottom.}
    \label{fig:environment}
\end{figure}


The environment also contains an overhead camera that sends images to the ROS system, where they are picked up by the computer vision nodes. This view only contains the apple, and can be seen in Figure \ref{fig:overhead}.

\begin{figure}[h]
    \centering
    \includegraphics[width=0.7\textwidth]{documents/latex/final_report/Overhead.png}
    \caption{View of the environment from the overhead camera.}
    \label{fig:overhead}
\end{figure}


Conceptually, the robot is treated as separate from the environment so that it can be swapped out, but it is positioned at the center of the scene, where it can reach both the groceries and the shelf while still facing a non-trivial placement task.


\subsubsection{Robot Arm and Gripper}

For this project, we used the robotic arm...
It had the advantage of being...

For the gripper, we used two blocks that open and close. While this designed only makes sense in the simulation, it's simplicity let us focus on other aspects of the work.

\subsubsection{ROS}

A list of the ROS nodes...

\subsection{Methodology}

For our experiment design we started off with a single apple object...


% for final report
\section{Experimental Results / Technical Demonstration}

A description of how you evaluated or demonstrated your solution.\footnote{a video of the robot doing x y z is available at...} 

% for proposal
%\section{Experiments and Evaluation}

%Describe how you will evaluate your approach/solution. What constitutes success? What metrics will you use? Do you have any preliminary hypothesis that you plan to test? Also, describe the RL domain or environment you plan to use. 

% for final report
\section{Conclusion and Future Work}

A high level summary of what was accomplished, along with a discussion on limitations and avenues for future work (typically 2 to 3 paragraphs). 


%\section{Timeline and Individual Responsibilities}

%State the timeline in terms of weeks and milestones you want to achieve. If working on a team, state what the individual responsibilities are at this point (i.e., who is going to do what, these may of course change over the course of the project). \cite{short2010no}.

\bibliographystyle{plain}
\bibliography{references}
\end{document}
